\documentclass[a4paper]{article}
\input{../../../../newpreamble.tex}
\title{210B: Lectures 8-10}
\author{Lance Remigio}
\begin{document}
\maketitle

\section*{Lecture 8}

\begin{prop}
    Every nonzero prime ideal in a PID is maximal.
\end{prop}
\vspace{5in}
\newpage
\begin{prop}
    If \( R  \) is any commutative ring such that the polynomial ring \( R[x]  \) is a PID, then \( R \) is a field.
\end{prop}
\newpage
\vspace{5in}
\begin{prop}
   In an integral domain, a prime element is irreducible.   
\end{prop}
\vspace{5in}
\newpage
\section*{Lecture 9}
\begin{prop}
   In a PID, a nonzero element is a prime if and only if it is irreducible. 
\end{prop}
\vspace{5in}
\newpage
\begin{prop}
   In a unit factorization domain, a non-zero element is prime if and only if it irreducible. 
\end{prop}

\vspace{5in}
\newpage
\section*{Lecture 10}
\begin{prop}
    Let \( I  \) be an ideal of \( R  \) and let \( (I) = I[x] \) (polynomials with coefficient in the ideal) generated by \( I  \) in \( R[x] \). (polynomials with coefficients in the ideal \( I \)) Then 
    \[  R[x] / (I) \cong (R /I ) [x]. \]
    In particular, if \( I  \) is a prime ideal of \( R  \), then \( (I) \) is a prime ideal.
\end{prop} 
\vspace{5in}
\newpage
\begin{theorem}[ ]
    Let \( F  \) be a field. Then \( F[x] \) is a Euclidean Domain. Specifically, if \( a(x), b(x) \in F[x] \) such that \( b(x) \neq 0  \). Then there exists a unique \( q(x), r(x) \in F[x] \) such that  
    \[  a(x) = q(x) b(x) + r(x) \]
    where \( r(x) = 0  \) or \( \text{deg}(r(x)) < \text{deg}(b(x))    \).
\end{theorem}
\vspace{5in}


\vspace{5in}

\end{document}

